\uuid{WPji}
\titre{Variables de loi normale centrée réduite}
\auteur{Maxime Nguyen}
\organisation{AMSCC}
\datecreate{2026-11-27}
\niveau{L2}
\module{Probabilité}
\chapitre{Probabilité continue}
\sousChapitre{Loi normale}
\difficulte{1}
\theme{loi normale, standardisation, symétrie}

\contenu{
\begin{question}
On note $\mathcal{N}(\mu,\sigma)$ la loi normale d'espérance $\mu$ et d'écart-type $\sigma$. Soit $X \sim \mathcal{N}(\mu, \sigma)$ avec $\mu \neq 0$ et $\sigma \neq 1$. Quelles variables suivent la loi $\mathcal{N}(0,1)$ ?
\begin{answers}
\good{$\dfrac{X - \mu}{\sigma}$}
\bad{$\dfrac{X - \mu}{\sigma^2}$}
\bad{$\dfrac{X}{\sigma} - \mu$}
\bad{$\left(\dfrac{X - \mu}{\sigma}\right)^2$}
\good{$\dfrac{\mu - X}{\sigma}$}
\end{answers}
\begin{explanations}
A est la standardisation du cours.

B a pour écart-type $\frac{\sigma}{\sigma^2} = \frac{1}{\sigma} \neq 1$.

C a pour espérance $\frac{\mu}{\sigma} - \mu$, qui n'est nulle que si $\sigma = 1$ : il faut centrer avant de diviser.

D est positive : elle ne peut pas suivre une loi normale (sa loi est la loi du $\chi^2$ à un degré de liberté).

E $= -\frac{X - \mu}{\sigma}$ ; si $Z \sim \mathcal{N}(0,1)$, alors $-Z \sim \mathcal{N}(0,1)$ par symétrie de la densité.

Contre-sens visés : diviser par la variance (B) ; diviser puis centrer avec
la mauvaise constante (C) ; refuser E parce que « le signe change ».
\end{explanations}
\end{question}
}
